\section{Ajustes al guion}
\label{sec:ajustes-guion}

Se revisaron las trece preguntas identificables en la transcripción del ensayo de entrevista estructurada con una participante que representó el rol de cajera de un establecimiento de productos eléctricos. La revisión consideró neutralidad, claridad, singularidad de la pregunta y adecuación al propósito de elicitación de requisitos \cite{pohl2025,ireb2026}. Se contrastó la formulación expresada durante el ensayo con el guion preparado antes de su aplicación.

De las trece preguntas examinadas, doce se consideraron adecuadas para su objetivo. No se identificaron preguntas inductivas, dobles, innecesariamente cerradas, con jerga técnica incomprensible ni hipotéticas sin necesidad. La primera pregunta, aunque cerrada, fue necesaria para solicitar la autorización de grabación. Se identificó una imprecisión en la formulación de la pregunta 9, que se detalla en la Tabla~\ref{tab:ajustes-guion}.

\begingroup
\footnotesize
\renewcommand{\arraystretch}{1.24}
\setlength{\LTpre}{8pt}
\setlength{\LTpost}{8pt}
\begin{longtable}{@{}p{0.045\textwidth}@{\hspace{5pt}}p{0.235\textwidth}@{\hspace{5pt}}p{0.14\textwidth}@{\hspace{5pt}}p{0.245\textwidth}@{\hspace{5pt}}p{0.245\textwidth}@{}}
\caption{Pregunta con defecto de formulación y versión corregida}\label{tab:ajustes-guion}\\
\toprule
\textbf{N.º} & \textbf{Pregunta tal como se formuló} & \textbf{Defecto} & \textbf{Explicación} & \textbf{Versión corregida}\\
\midrule
\endfirsthead
\multicolumn{5}{@{}l}{\textit{Tabla \thetable\ (continuación)}}\\
\toprule
\textbf{N.º} & \textbf{Pregunta tal como se formuló} & \textbf{Defecto} & \textbf{Explicación} & \textbf{Versión corregida}\\
\midrule
\endhead
\bottomrule
\endfoot
9 & Si ha ocurrido un error al registrar una venta, ¿cómo lo corrigieron? & Ambigüedad leve & La pregunta se dirige a una persona en singular mediante ``ha ocurrido'', pero cambia a la tercera persona del plural en ``corrigieron'', sin identificar quién realizó la corrección. Esto puede generar dudas sobre si se solicita el procedimiento personal o el del establecimiento. & ¿Cómo se corrige un error cuando se detecta en el registro de una venta?\\
\end{longtable}
\endgroup

La reformulación mantiene la finalidad de conocer el procedimiento de corrección de errores, conserva el carácter abierto y elimina la referencia imprecisa al responsable. Además, permite describir las diferencias entre una corrección anterior a la autorización del comprobante y una corrección posterior, distinción que apareció en la respuesta del ensayo. La nueva pregunta no presupone el método de corrección ni sugiere una solución informática.

En la pregunta 7 se observó una variación de redacción respecto del guion preparado: durante el ensayo se empleó ``¿Qué paso de la facturación le toma más tiempo, si es que lo hay?''. Esta variación conserva una pregunta abierta y neutral, por lo que no se clasificó como defecto metodológico. Las restantes preguntas mantienen un propósito claro y no requieren correcciones en esta revisión.

Los ajustes se refieren a la calidad del guion empleado en una simulación académica. Las respuestas del juego de roles aportan información para discutir necesidades potenciales, pero no constituyen evidencia directa obtenida de trabajadores reales del sector comercial.
